\documentclass[11pt,a4paper]{amsart}

\usepackage[T1]{fontenc}
\usepackage{lmodern}
\usepackage{amsmath,amssymb,mathtools}
\usepackage[protrusion=true,expansion=false]{microtype}
\usepackage[a4paper,margin=30mm]{geometry}
\usepackage[dvipsnames]{xcolor}
\usepackage[colorlinks=true,linkcolor=MidnightBlue,citecolor=MidnightBlue,urlcolor=MidnightBlue]{hyperref}
\hypersetup{
  pdftitle={Subpower Energy and Dyadic Defect Criteria for the Riemann Hypothesis},
  pdfauthor={Mayk Loide Baccaro}
}

\newtheorem{theorem}{Theorem}[section]
\newtheorem{proposition}[theorem]{Proposition}
\newtheorem{lemma}[theorem]{Lemma}
\newtheorem{corollary}[theorem]{Corollary}
\theoremstyle{definition}
\newtheorem{definition}[theorem]{Definition}
\theoremstyle{remark}
\newtheorem{remark}[theorem]{Remark}

\newcommand{\C}{\mathbb C}
\newcommand{\R}{\mathbb R}
\newcommand{\Pone}{\mathrm{P1}}
\setlength{\emergencystretch}{2em}

\title[Energy and dyadic defect criteria]{Subpower Energy and Dyadic Defect Criteria for the Riemann Hypothesis}
\author[Mayk Loide Baccaro]{Mayk Loide Baccaro\\
Independent researcher\\
\href{mailto:m.baccaro7663@student.nu.edu}{\texttt{m.baccaro7663@student.nu.edu}}}
\date{20 August 2026}
\subjclass[2020]{Primary 11M26; Secondary 11N37, 11M06}
\keywords{Riemann hypothesis, M\"obius function, Mellin transform, dyadic decomposition, reciprocal zeta function}

\begin{document}

\begin{abstract}
Let
\[
 b_n=-\mu(n)\log n,
 \qquad B_N=\sum_{n\leq N}b_n.
\]
We prove that the Riemann hypothesis is equivalent to subpower growth of the
positive prefix energy
\[
 \sum_{n\leq X}\frac{B_n^2}{n(n+1)},
\]
and also to convergence of each positive power-damped version of this sum.
The same criterion has a continuous form: for every $\sigma>1/2$, the
weighted energy of the right-continuous step function $B(x)=B_{\lfloor
x\rfloor}$ is finite, or equivalently its Mellin transform is square
integrable on $\Re s=\sigma$.  An exact prefix-square identity then gives
equivalent cumulative and individual lower bounds for signed dyadic defects.
As a stronger sufficient condition, eventual nonnegativity of the dyadic
defects implies the Riemann hypothesis, simplicity of all nontrivial zeros,
and the explicit bound
\[
 \sum_{\rho}\left|\frac{1}{\rho\zeta'(\rho)}\right|^2
 \leq \frac{18}{\pi^2}.
\]
The paper proves these equivalences, but it does not establish any of the
equivalent estimates or the eventual-sign hypothesis.  It therefore does
not prove the Riemann hypothesis.
\end{abstract}

\maketitle

\section{Introduction}

The classical Littlewood criterion expresses the Riemann hypothesis through
cancellation in the summatory M\"obius function.  This paper studies the
closely related coefficients
\[
 b_n=-\mu(n)\log n=(\mu*\Lambda)(n),
\]
which are the Dirichlet coefficients of the derivative of $1/\zeta(s)$.
Their prefix sums $B_N$ lead to two views of the same criterion.

The first view is positive.  It asks whether the cumulative square energy
\[
 A(X)=\sum_{2\leq n\leq X}\frac{B_n^2}{n(n+1)}
\]
grows more slowly than every fixed positive power of $X$.  Positivity makes
this condition stable under damping and lets it pass directly to a weighted
Mellin $L^2$ norm.

The second view is signed and local.  On the block $(L,R]$, define
\[
 D(L,R)=\sum_{L<n\leq R}
 \frac{b_n^2-b_nB_{n-1}}{n}.
\]
An exact identity relates $D(L,R)$ to the positive prefix energy.  On the
power-of-two blocks $(2^j,2^{j+1}]$, this identity shows that subpower lower
bounds for either each defect or their cumulative sum are equivalent to the
Riemann hypothesis.

The principal theorem is an equivalence theorem, not an estimate for the
new quantities.  It does not prove any of its criteria.  A separate corollary
starts from the stronger, also unproved, hypothesis that the power-of-two
defects are eventually nonnegative.  That sign condition supplies a sharp
leading constant in the prefix energy.  Mellin-Plancherel then yields an
explicit boundary norm.  The poles of the differentiated reciprocal zeta
function consume nonnegative portions of this norm, which forces simple
zeros and the stated reciprocal-residue budget.

Several $L^2$ formulations of the Riemann hypothesis are already known.
The Nyman-Beurling theory and its refinements use approximation or closure of
families of fractional-part functions
\cite{Nyman1950,Beurling1955,BaezDuarte2003,Kapustin2019}.  The weak Mertens
theory uses mean squares of the ordinary Mertens function and has strong
consequences for reciprocal zeta residues \cite{Titchmarsh1986,Ng2004}.
Modified mollifiers, M\"obius-convolution criteria, generalized von Mangoldt
sums, and smoothed M\"obius transforms provide further nearby formulations
\cite{ConreyMyerson2000,BaezDuarte2005,BanksSinha2022,Verjovsky2026}.
The criterion here instead uses the literal prefix of $-\mu(n)\log n$, its
positive weighted square energy, and an exact signed dyadic identity.  These
comparisons locate the argument; they do not establish a novelty or priority
claim.

\section{The exact block identity}

Let $\mu$ be the M\"obius function and $\Lambda$ the von Mangoldt function.
For $n\geq1$, set
\begin{equation}\label{eq:bB}
 b_n=(\mu*\Lambda)(n)=-\mu(n)\log n,
 \qquad
 B_N=\sum_{n\leq N}b_n.
\end{equation}
The coefficient identity follows by differentiating
$1/\zeta(s)=\sum_{n\geq1}\mu(n)n^{-s}$ in $\Re s>1$.  We also use the
right-continuous step function
\begin{equation}\label{eq:step}
 B(x)=\sum_{n\leq x}b_n=B_{\lfloor x\rfloor},
 \qquad x\geq1.
\end{equation}
Values at the integer endpoints do not affect the integrals below, but this
convention fixes the discrete-to-continuous correspondence.

For integers $2\leq L<R$, define
\begin{align}
 D(L,R)&=\sum_{L<n\leq R}
 \frac{b_n^2-b_nB_{n-1}}{n},\label{eq:D}\\
 V(L,R)&=\sum_{L<n\leq R}
 \frac{\mu(n)^2\log^2 n}{n}.\label{eq:V}
\end{align}

\begin{proposition}[Prefix-square identity]\label{prop:identity}
For every $2\leq L<R$,
\begin{equation}\label{eq:identity}
 2D(L,R)=3V(L,R)-H(L,R),
\end{equation}
where
\begin{equation}\label{eq:H}
 H(L,R)=\frac{B_R^2}{R}-\frac{B_L^2}{L+1}
 +\sum_{L<n<R}\frac{B_n^2}{n(n+1)}.
\end{equation}
\end{proposition}

\begin{proof}
Since $b_n=B_n-B_{n-1}$,
\[
 2b_nB_{n-1}=B_n^2-B_{n-1}^2-b_n^2.
\]
It follows that
\[
 2(b_n^2-b_nB_{n-1})
 =3b_n^2-B_n^2+B_{n-1}^2.
\]
After division by $n$, summation over $L<n\leq R$ gives
\[
 \sum_{L<n\leq R}\frac{B_n^2-B_{n-1}^2}{n}
 =\frac{B_R^2}{R}-\frac{B_L^2}{L+1}
 +\sum_{L<n<R}\frac{B_n^2}{n(n+1)}.
\]
Now use $b_n^2=\mu(n)^2\log^2 n$.
\end{proof}

For $j\geq1$, write
\begin{equation}\label{eq:dyadic-defects}
 D_j=D(2^j,2^{j+1}),
 \qquad
 S_J=\sum_{j=1}^{J-1}D_j\quad(J\geq2).
\end{equation}
Summing Proposition~\ref{prop:identity} over adjacent dyadic blocks gives
the following exact global form.

\begin{corollary}[Exact dyadic balance]\label{cor:dyadic-identity}
For $X=2^J$ with $J\geq2$,
\begin{equation}\label{eq:dyadic-identity}
 2S_J+A(X)+\frac{B_X^2}{X+1}
 =3W(X)+\frac{B_2^2}{2},
\end{equation}
where
\[
 W(X)=\sum_{2<n\leq X}\frac{\mu(n)^2\log^2n}{n}.
\]
\end{corollary}

\begin{proof}
At an interior dyadic boundary $n$, the adjacent endpoint terms combine as
\[
 \frac{B_n^2}{n}-\frac{B_n^2}{n+1}
 =\frac{B_n^2}{n(n+1)}.
\]
\[
 \sum_{j=1}^{J-1}H(2^j,2^{j+1})
 =\frac{B_X^2}{X}-\frac{B_2^2}{3}
 +\sum_{2<n<X}\frac{B_n^2}{n(n+1)}.
\]
Substitution in Proposition~\ref{prop:identity}, followed by
\[
 A(X)=\frac{B_2^2}{6}
 +\sum_{2<n<X}\frac{B_n^2}{n(n+1)}
 +\frac{B_X^2}{X(X+1)},
\]
gives \eqref{eq:dyadic-identity}.
\end{proof}

We shall use the standard squarefree estimate and its partial-summation
consequence
\begin{align}
 \sum_{n\leq X}\mu(n)^2&=\frac{6}{\pi^2}X+O(\sqrt X),\label{eq:squarefree}\\
 \sum_{n\leq X}\frac{\mu(n)^2\log^2n}{n}
 &=\frac{2}{\pi^2}\log^3X+O(\log^2X).\label{eq:weighted-squarefree}
\end{align}

\section{Positive energy, Mellin energy, and dyadic defects}

Define the positive cumulative energy
\begin{equation}\label{eq:A}
 A(X)=\sum_{2\leq n\leq X}\frac{B_n^2}{n(n+1)}.
\end{equation}
For $\eta>0$ and $\sigma>1/2$, also define
\begin{align}
 \mathcal E_\eta
 &=\sum_{n\geq2}\frac{B_n^2}{n(n+1)}n^{-\eta},\label{eq:weighted-energy}\\
 I(\sigma)
 &=\int_1^\infty |B(x)|^2x^{-2\sigma-1}\,dx.\label{eq:I}
\end{align}

\begin{lemma}[Exact continuous realization]\label{lem:mellin}
For every $\sigma>1/2$,
\begin{equation}\label{eq:step-sum}
 I(\sigma)=\sum_{n\geq2}B_n^2
 \frac{n^{-2\sigma}-(n+1)^{-2\sigma}}{2\sigma}.
\end{equation}
If $\eta=2\sigma-1$, then
\begin{equation}\label{eq:comparison}
 I(\sigma)<\infty
 \quad\Longleftrightarrow\quad
 \mathcal E_\eta<\infty.
\end{equation}
Put
\[
 g_\sigma(y)=B(e^y)e^{-\sigma y}\mathbf 1_{[0,\infty)}(y)
\]
and, for $T>0$, let
\[
 F_{\sigma,T}(t)=\int_0^T g_\sigma(y)e^{-ity}\,dy.
\]
Then $I(\sigma)<\infty$ if and only if $F_{\sigma,T}$ converges in
$L^2(\R)$ as $T\to\infty$.  Denote the limit by $F_\sigma$.  In that case,
\begin{equation}\label{eq:plancherel}
 \int_{-\infty}^{\infty}|F_\sigma(t)|^2\,dt
 =2\pi I(\sigma).
\end{equation}
\end{lemma}

\begin{proof}
By \eqref{eq:step}, $B(x)=B_n$ on $[n,n+1)$ apart from an irrelevant
endpoint.  Integration on each interval gives \eqref{eq:step-sum}.  The
mean-value theorem gives
\[
 \frac{n^{-2\sigma}-(n+1)^{-2\sigma}}{2\sigma}
 \asymp_\sigma n^{-2\sigma-1}
 =n^{-2-\eta}
 \asymp \frac{n^{-\eta}}{n(n+1)},
\]
which proves \eqref{eq:comparison}.  The substitution $x=e^y$ gives
\[
 \int_0^\infty|g_\sigma(y)|^2\,dy=I(\sigma).
\]
For $T>U$, Plancherel for the convention
$\widehat g(t)=\int_\R g(y)e^{-ity}\,dy$ gives
\[
 \|F_{\sigma,T}-F_{\sigma,U}\|_2^2
 =2\pi\int_U^T|g_\sigma(y)|^2\,dy.
\]
Thus $F_{\sigma,T}$ is Cauchy in $L^2$ exactly when $I(\sigma)$ is finite.
Taking the limit proves \eqref{eq:plancherel}.
\end{proof}

\begin{theorem}[Equivalent subpower criteria]\label{thm:equivalence}
The following statements are equivalent.
\begin{enumerate}
 \item The Riemann hypothesis is true.
 \item For every $\varepsilon>0$,
 \begin{equation}\label{eq:PE}
  A(X)=O_\varepsilon(X^\varepsilon).
 \end{equation}
 \item For every $\eta>0$, the series $\mathcal E_\eta$ converges.
 \item For every $\sigma>1/2$, the integral $I(\sigma)$ is finite.
 \item For every $\sigma>1/2$, the truncated transforms
 $F_{\sigma,T}$ in Lemma~\ref{lem:mellin} converge in $L^2(\R)$.
 \item For every $\varepsilon>0$, there is $C_\varepsilon$ such that
 \begin{equation}\label{eq:CDD}
  S_J\geq-C_\varepsilon 2^{\varepsilon J}
  \qquad(J\geq2).
 \end{equation}
 \item For every $\varepsilon>0$, there is $C_\varepsilon$ such that
 \begin{equation}\label{eq:IDD}
  D_j\geq-C_\varepsilon 2^{\varepsilon j}
  \qquad(j\geq1).
 \end{equation}
\end{enumerate}
\end{theorem}

\begin{proof}
We first connect the positive discrete conditions.  If
\eqref{eq:PE} holds and $\eta>0$, use the bound with exponent $\eta/2$
and partial summation against the decreasing weight $x^{-\eta}$; this
proves $\mathcal E_\eta<\infty$.  Conversely, positivity gives
\[
 A(X)\leq X^\eta\mathcal E_\eta.
\]
Taking $\eta=\varepsilon$ proves \eqref{eq:PE}.  Lemma~\ref{lem:mellin}
then proves the equivalence of statements 2 through 5.

We next show that statement 3 implies RH.  On $\Re s>1$, partial
summation gives
\begin{equation}\label{eq:dirichlet-mellin}
 F(s):=\int_1^\infty B(x)x^{-s-1}\,dx
 =\frac{1}{s}\left(\frac{1}{\zeta(s)}\right)'.
\end{equation}
Let $K$ be a compact subset of $\Re s>1/2$, and choose
$\sigma_0$ with $1/2<\sigma_0<\min_{s\in K}\Re s$.  Statement 3 and
Lemma~\ref{lem:mellin} give $I(\sigma_0)<\infty$.  Cauchy-Schwarz yields
\[
 \int_1^\infty |B(x)|x^{-\Re s-1}\,dx
 \leq I(\sigma_0)^{1/2}
 \left(\int_1^\infty
 x^{-2(\Re s-\sigma_0)-1}\,dx\right)^{1/2},
\]
uniformly on $K$.  Thus the integral defining $F$ is locally uniformly
convergent and holomorphic in $\Re s>1/2$.

The half-plane is simply connected, so $sF(s)$ has a holomorphic primitive
$Q(s)$ there.  Normalize it at one point in $\Re s>1$ to equal
$1/\zeta(s)$.  In $\Re s>1$, equation \eqref{eq:dirichlet-mellin} gives
$Q'=(1/\zeta)'$, so the normalization gives $Q=1/\zeta$.  On
$\{\Re s>1/2\}\setminus\{1\}$, the meromorphic identity theorem gives
\[
 \zeta(s)Q(s)=1.
\]
Hence $\zeta$ has no zero with real part greater than $1/2$.  The
functional equation reflects every nontrivial zero across the critical
line, so all nontrivial zeros lie on that line.  This proves RH.
Moreover, $Q'=sF$ and $Q=1/\zeta$ give
\[
 F(s)=\frac{1}{s}\left(\frac{1}{\zeta(s)}\right)'
 \qquad(\Re s>1/2,\ s\neq1).
\]
Both sides extend holomorphically across $s=1$.
The same meromorphic expression therefore controls the approach from the
right to a zero on the critical line.

Conversely, RH implies the classical estimate
\cite[Theorem~14.25(C)]{Titchmarsh1986}
\[
 M(x):=\sum_{n\leq x}\mu(n)
 =O_\alpha(x^{1/2+\alpha})
\]
for every $\alpha>0$.  Partial summation in \eqref{eq:bB} gives
\begin{equation}\label{eq:B-from-M}
 B(x)=-M(x)\log x+\int_1^x\frac{M(t)}{t}\,dt.
\end{equation}
After absorbing the logarithm into a slightly larger exponent,
$B(x)=O_\alpha(x^{1/2+\alpha})$ for every $\alpha>0$.  Substitution in
\eqref{eq:A} proves \eqref{eq:PE}.

It remains to connect the dyadic statements.  Statement 7 implies statement
6 by summing a geometric progression.  If statement 6 holds, then
Corollary~\ref{cor:dyadic-identity} and \eqref{eq:weighted-squarefree} give,
for $X=2^J$,
\[
 A(X)\leq 3W(X)+2C_\varepsilon X^\varepsilon+O(1)
 =O_\varepsilon(X^\varepsilon).
\]
Monotonicity of $A$ extends this bound to every $X$, so statement 2 follows.

Finally assume RH.  Choose $\alpha>0$ so small that $2\alpha<\varepsilon$.
The bound after \eqref{eq:B-from-M} and \eqref{eq:H} give
\[
 H(2^j,2^{j+1})=O_\alpha(2^{2\alpha j}).
\]
Since $V\geq0$, Proposition~\ref{prop:identity} yields
\[
 D_j\geq-O_\alpha(2^{2\alpha j}),
\]
which is statement 7.  All seven statements are therefore equivalent.
\end{proof}

\begin{remark}[Boundary not included]\label{rem:boundary}
The Mellin conditions in Theorem~\ref{thm:equivalence} concern every fixed
line $\Re s=\sigma>1/2$.  The theorem does not assert square integrability
on the critical line itself.
\end{remark}

\section{The stronger eventual-sign condition}

The equivalence theorem allows small negative dyadic defects of subpower
size.  Eventual nonnegativity is stronger and supplies an explicit leading
constant.

\begin{theorem}[Eventual sign and boundary budget]\label{thm:eventual}
Suppose there is $j_0$ such that
\begin{equation}\label{eq:eventual}
 D_j\geq0\qquad(j\geq j_0).
\end{equation}
Then the Riemann hypothesis is true, every nontrivial zero of $\zeta$ is
simple, and
\begin{equation}\label{eq:budget}
 \sum_\rho\left|\frac{1}{\rho\zeta'(\rho)}\right|^2
 \leq\frac{18}{\pi^2}.
\end{equation}
The sum includes both conjugate ordinates.  In addition,
\begin{equation}\label{eq:A-sharp}
 A(X)\leq\frac{6}{\pi^2}\log^3X+O(\log^2X)
\end{equation}
and, with $F$ as in \eqref{eq:dirichlet-mellin},
\begin{equation}\label{eq:hardy-ceiling}
 \limsup_{\delta\downarrow0}\delta^3
 \int_{-\infty}^{\infty}
 \left|F\left(\frac12+\delta+it\right)\right|^2\,dt
 \leq\frac{9}{\pi}.
\end{equation}
\end{theorem}

\begin{proof}
Condition \eqref{eq:eventual} and
Corollary~\ref{cor:dyadic-identity} give
\[
 A(2^J)\leq3W(2^J)+O(1).
\]
Equation \eqref{eq:weighted-squarefree} proves \eqref{eq:A-sharp} first at
powers of two and then, by monotonicity, for all $X$.
Theorem~\ref{thm:equivalence} now proves RH.

For $\delta>0$, the step convention gives
\begin{align*}
 I\left(\frac12+\delta\right)
 &=\sum_{n\geq2}B_n^2\int_n^{n+1}x^{-2-2\delta}\,dx\\
 &\leq\sum_{n\geq2}
 \frac{B_n^2}{n(n+1)}n^{-2\delta}
 =\int_{2^-}^{\infty}x^{-2\delta}\,dA(x).
\end{align*}
Integration by parts and \eqref{eq:A-sharp} show that
\begin{equation}\label{eq:I-ceiling}
 \limsup_{\delta\downarrow0}\delta^3
 I\left(\frac12+\delta\right)
 \leq\frac{9}{2\pi^2}.
\end{equation}
Indeed,
\[
 \frac{6}{\pi^2}\,2\delta
 \int_1^\infty \log^3x\,x^{-2\delta-1}\,dx
 =\frac{9}{2\pi^2\delta^3},
\]
while the $O(\log^2x)$ term contributes $O(\delta^{-2})$.
Equation \eqref{eq:plancherel} turns \eqref{eq:I-ceiling} into
\eqref{eq:hardy-ceiling}.

Let $\rho=1/2+i\gamma$ be a nontrivial zero of multiplicity $m$, and write
\[
 \zeta(s)=(s-\rho)^m g(s),\qquad g(\rho)\neq0.
\]
In a right-hand neighborhood of $\rho$, the preceding meromorphic identity
gives
\begin{equation}\label{eq:laurent-general}
 F(s)=c_\rho(s-\rho)^{-m-1}
 +O_\rho(|s-\rho|^{-m}),
 \qquad c_\rho=-\frac{m}{\rho g(\rho)}.
\end{equation}
Choose $\eta>0$ so that the Laurent estimate holds in the right half of
$|s-\rho|<2\eta$, the disk avoids $0$ and $1$, and $\rho$ is its only
zeta zero.  Put $I_\rho=[\gamma-\eta,\gamma+\eta]$.  Scaling
$t-\gamma=\delta u$ in \eqref{eq:laurent-general} gives
\begin{equation}\label{eq:local-energy-general}
 \int_{I_\rho}|F(1/2+\delta+it)|^2\,dt
 =|c_\rho|^2C_m\delta^{-2m-1}+o_\rho(\delta^{-2m-1}),
\end{equation}
where
\[
 C_m=\int_{-\infty}^{\infty}(1+u^2)^{-m-1}\,du>0.
\]
Indeed, the squared remainder in \eqref{eq:laurent-general} contributes
$O_\rho(\delta^{1-2m})$, and the cross term contributes
$O_\rho(\delta^{-2m})$.  If $m\geq2$, multiplying
\eqref{eq:local-energy-general} by $\delta^3$ makes its left side tend to
infinity, contradicting the finite global limsup in
\eqref{eq:hardy-ceiling}.  Thus every nontrivial zero is simple.

For $m=1$, $g(\rho)=\zeta'(\rho)$ and
\[
 C_1=\int_{-\infty}^{\infty}\frac{du}{(1+u^2)^2}=\frac{\pi}{2}.
\]
Consequently, \eqref{eq:local-energy-general} becomes
\begin{equation}\label{eq:local-energy}
 \int_{I_\rho}|F(1/2+\delta+it)|^2\,dt
 =\frac{\pi}{2\delta^3}
 \left|\frac{1}{\rho\zeta'(\rho)}\right|^2
 +o_\rho(\delta^{-3}).
\end{equation}
Fix a finite set $E$ of zeros and choose the intervals $I_\rho$ pairwise
disjoint.  The global norm is nonnegative.  Summing
\eqref{eq:local-energy} over $E$ and using \eqref{eq:hardy-ceiling} gives
\[
 \frac{\pi}{2}\sum_{\rho\in E}
\left|\frac{1}{\rho\zeta'(\rho)}\right|^2
 \leq\frac{9}{\pi}.
\]
Taking the supremum over all finite $E$ proves \eqref{eq:budget}.
\end{proof}

\section{Relation to the P1 condition}

A previously considered condition asks for more than
\eqref{eq:eventual}.  For every integer $L\geq2$ and
$R\in\{2L,2L+1\}$, let $D(L,R)$ be as in \eqref{eq:D}.

\begin{definition}[P1]\label{def:P1}
Condition $\Pone$ asserts
\[
 D(L,R)\geq0
\]
for every $L\geq2$ and both permitted endpoints $R=2L$ and $R=2L+1$.
\end{definition}

Taking $L=2^j$ shows that $\Pone$ implies \eqref{eq:eventual}, with
$j_0=1$.  Theorem~\ref{thm:eventual} therefore gives
\[
 \Pone\quad\Longrightarrow\quad
 \text{RH, simple nontrivial zeros, and \eqref{eq:budget}}.
\]
This is only a one-way implication.  Theorem~\ref{thm:equivalence} proves
that the subpower energy and defect bounds are equivalent to RH, but it does
not prove that RH implies $\Pone$.  It also does not prove a strict logical
separation between $\Pone$ and RH.  Condition $\Pone$ remains unproved, and
finite verification of its inequalities cannot establish the universal
statement.

\section{Context and limitations}

The Mellin map in Lemma~\ref{lem:mellin} places the positive energy criterion
in a classical analytic setting, but its role differs from the
Nyman-Beurling program.  Nyman and Beurling characterize RH by density of
specified fractional-part functions in an $L^2$ space
\cite{Nyman1950,Beurling1955}.  B\'aez-Duarte restricts the admissible
generators to an arithmetic family \cite{BaezDuarte2003}.  Kapustin gives
unitarily equivalent Hardy-space and weighted-$L^2$ approximation models,
including a closure criterion for the span of $\rho(nx)$
\cite[Theorem~1]{Kapustin2019}.  Here there is no approximation family: the tested
object is the single step function formed by the literal prefixes of
$-\mu(n)\log n$, and the condition is finiteness of its weighted square
energy on every line to the right of the critical line.

The weak Mertens condition also uses a positive mean square and implies RH,
simple zeros, and summability of reciprocal residues
\cite{Titchmarsh1986,Ng2004}.  Its summatory function is
$M(x)=\sum_{n\leq x}\mu(n)$, whereas the present energy uses the logarithmic
prefix $B(x)$ and arises from $(1/\zeta)'$.  The explicit constant in
Theorem~\ref{thm:eventual} comes from the exact dyadic identity and the
leading squarefree density.

The criteria of Conrey and Myerson use modified M\"obius mollifiers in the
Balazard-Saias framework \cite{ConreyMyerson2000}; B\'aez-Duarte develops
general M\"obius-convolution criteria \cite{BaezDuarte2005}; Banks and Sinha
use uniform twisted sums of generalized von Mangoldt functions
\cite{BanksSinha2022}; and Verjovsky studies exponential smoothing and local
Fourier formulations of M\"obius cancellation \cite{Verjovsky2026}.  These
works share coefficients or transform methods with the present paper, but
their stated criteria are different.

No computation is used in any proof above.  In particular, finite tests of
the signed defects and finite partial sums of \eqref{eq:budget} cannot prove
the universal hypotheses or control their tails.  The paper establishes
equivalences and conditional consequences only.

\begin{thebibliography}{99}

\bibitem{BaezDuarte2003}
L.~B\'aez-Duarte,
\emph{A strengthening of the Nyman-Beurling criterion for the Riemann hypothesis},
Atti Accad. Naz. Lincei Cl. Sci. Fis. Mat. Natur. Rend. Lincei (9)
Mat. Appl. \textbf{14} (2003), no.~1, 5--11,
\href{https://arxiv.org/abs/math/0202141}{arXiv:math/0202141}.

\bibitem{BaezDuarte2005}
L.~B\'aez-Duarte,
\emph{M\"obius-convolutions and the Riemann hypothesis},
Ramanujan J. \textbf{9} (2005), 11--23,
\href{https://doi.org/10.1007/s11139-005-0839-0}{doi:10.1007/s11139-005-0839-0}.

\bibitem{BanksSinha2022}
W.~D. Banks and S.~Sinha,
\emph{The Riemann hypothesis via the generalized von Mangoldt function},
arXiv:2209.11768v2,
\href{https://arxiv.org/abs/2209.11768}{arXiv:2209.11768}.

\bibitem{Beurling1955}
A.~Beurling,
\emph{A closure problem related to the Riemann zeta-function},
Proc. Natl. Acad. Sci. USA \textbf{41} (1955), 312--314,
\href{https://doi.org/10.1073/pnas.41.5.312}{doi:10.1073/pnas.41.5.312}.

\bibitem{ConreyMyerson2000}
J.~B. Conrey and G.~Myerson,
\emph{On the Balazard-Saias criterion for the Riemann hypothesis},
arXiv:math/0002254,
\href{https://arxiv.org/abs/math/0002254}{arXiv:math/0002254}.

\bibitem{Kapustin2019}
V.~V. Kapustin,
\emph{Beurling's theorem, Davenport's formula, and the Riemann hypothesis},
St. Petersburg Math. J. \textbf{30} (2019), no.~6, 917--932,
\href{https://doi.org/10.1090/spmj/1577}{doi:10.1090/spmj/1577}.

\bibitem{Ng2004}
N.~Ng,
\emph{The distribution of the summatory function of the M\"obius function},
Proc. London Math. Soc. (3) \textbf{89} (2004), 361--389,
\href{https://doi.org/10.1112/S0024611504014741}{doi:10.1112/S0024611504014741}.

\bibitem{Nyman1950}
B.~Nyman,
\emph{On some groups and semi-groups of translations},
Ph.D. thesis, Uppsala University, 1950.

\bibitem{Titchmarsh1986}
E.~C. Titchmarsh,
\emph{The Theory of the Riemann Zeta-Function},
2nd ed., revised by D.~R. Heath-Brown,
Oxford University Press, Oxford, 1986.

\bibitem{Verjovsky2026}
A.~Verjovsky,
\emph{How random is the M\"obius function? Smoothing, probability, and the Riemann hypothesis},
arXiv:2607.25002v2,
\href{https://arxiv.org/abs/2607.25002}{arXiv:2607.25002}.

\end{thebibliography}

\end{document}
